\subsection{Evaluation}
\label{sec:eval}

Two quantities are reported per episode and never combined. \emph{Realised regret} is read
from what the campaign did; the \emph{regret a perfect reasoner would still carry} is read
from what each design could have supported, before its outcome is seen. Neither elicits
anything from the agent beyond a design and a recommendation.

\subsubsection{Realised regret}
\label{sec:outcome}

Both quantities of \S\ref{sec:conditions} are made precise here. They are regret against
the same optimum in the same currency, because the experiment is paid for in the units of
what it is trying to maximise. The first is \eqref{eq:problem}, elicited after
\emph{every} round rather than only at $R$, and averaged over instances:
\begin{equation}
\bar{\mathcal{R}}(r) \;=\; \mathbb{E}_{\theta \sim P_\Theta}
\Big[\, f_\theta(x^\star_\theta) - f_\theta(\hat{x}_r) \,\Big],
\qquad x^\star_\theta \in \arg\max_{x \in \mathcal{X}} f_\theta(x).
\label{eq:regret}
\end{equation}
The second is what the experiments themselves forwent: every occupied slot is run at the
treatment it was given, so a slot assigned a poor one costs the gap to the optimum for a
full cycle of length $\Delta$,
\begin{equation}
\mathcal{R}_{\mathrm{cum}} \;=\; \Delta \sum_{r=1}^{R} \sum_{u \in \mathcal{U}_r}
\Big[\, f_\theta(x^\star_\theta) - f_\theta(x_u) \,\Big].
\label{eq:cumregret}
\end{equation}
This is the usual simple-versus-cumulative pair, and the second term is usually dropped in
benchmark settings because a trial there costs nothing. Here an agent that locates the optimum
by first running every bad treatment has answered the question and exhausted the facility
that paid for it. Neither
implies the other, and in our runs they do order agents differently
(Appendix~\ref{app:refs}); an agent that runs no experiments forgoes nothing and is judged
by \eqref{eq:regret} alone.

\subsubsection{Design and decision quality}
\label{sec:eig}

A campaign can fail in two ways that a single score cannot tell apart: the experiments could
not have answered the question, or they could and the agent did not read them. Both are
measured against $\mathcal{R}^\star(D_r)$, the regret still left after a decision-maker who
reasoned perfectly from this design's outcome,
\begin{equation}
\mathcal{R}^\star(D_r) \;=\;
\underbrace{\mathbb{E}_{p(y \mid D_r)}}_{\text{over outcomes not yet seen}}
\Big[\, \underbrace{\min_{x \in \mathcal{X}}}_{\text{best recommendation}} \;
\underbrace{\mathbb{E}_{p(\theta \mid y, D_r)}
\big[\, f_\theta(x^\star_\theta) - f_\theta(x) \,\big]}_{\text{its expected regret}} \Big].
\label{eq:eig}
\end{equation}
Read it from the inside out. The bracket is the regret \eqref{eq:problem} of recommending
$x$ at site $\theta$, averaged over the posterior a perfect reasoner would hold after seeing
$y$, the vector of observations \eqref{eq:obs} that design $D_r$ returns. The $\min$ picks
the recommendation that posterior supports best. The outer expectation averages over
outcomes the design has not produced yet, so $\mathcal{R}^\star$ is a property of the design
alone and can be computed before it is run. No agent can compute it: doing so needs
$P_\Theta$ and the forward model, which the interface withholds. Written
$\mathcal{R}^\star(\varnothing)$ for a design that runs nothing, it is the regret of the
best single recommendation, and so the most any experiment on the task could remove.

From it we form one ratio for what was run and one for what was concluded,
\begin{equation}
\underbrace{c(D_r) \;=\; \frac{\mathcal{R}^\star(\varnothing) - \mathcal{R}^\star(D_r)}
                  {\mathcal{R}^\star(\varnothing)}}_{\text{\emph{design efficiency}}},
\qquad
\underbrace{\eta(D_r) \;=\; \frac{\mathcal{R}^\star(D_r)}{\bar{\mathcal{R}}(r)}}
           _{\text{\emph{decision efficiency}}}.
\label{eq:capture}
\end{equation}
\emph{Design efficiency} $c$ is the fraction of the available risk
$\mathcal{R}^\star(\varnothing)$ that the experiments removed; it is $0\%$ for a design that
teaches nothing and $100\%$ for one that identifies the optimum exactly. \emph{Decision
efficiency} $\eta$ divides what the design made available, $\mathcal{R}^\star(D_r)$, by what
the agent actually achieved with it, $\bar{\mathcal{R}}(r)$ from \eqref{eq:regret}; it is
near $100\%$ when the recommendation is as good as the data allowed and near $0\%$ when the
agent left most of it unused. The \emph{decision} here is the recommendation $\hat{x}_r$
and nothing else: both terms of $\eta$ are the regret of a point in $\mathcal{X}$, against
the same optimum, and differ only in who chose it --- the $\min$ inside \eqref{eq:eig}, or
the agent. Both are reported as percentages, larger is better for both,
and an agent can fail at either.

Neither is new. $\mathcal{R}^\star(\varnothing)$ and $\mathcal{R}^\star(D)$ are the prior
and preposterior Bayes risk of the decision problem, so their difference is the expected
value of sample information \citep{raiffa1961,howard1966}, equivalently a Knowledge Gradient
\citep{frazier2008kg} for a batch rather than a single point. What is ours is the pairing.
Unlike information gain about $\theta$ \citep{boxinggym}, \eqref{eq:eig} scores the
recommendation rather than the parameters, so a design that pins $\theta$ down while
spending its budget far from the optimum scores low; and since \eqref{eq:eig} contains no
$\hat{x}_r$, $c$ is a functional of the design and the instance distribution alone while
$\eta$ carries $\bar{\mathcal{R}}(r)$ in its denominator. The two cannot be restatements of
each other, and \S\ref{sec:results} shows one intervention moving them in opposite
directions.

\paragraph{Estimation.}\label{sec:estimator} We discretise $P_\Theta$ into $M$ atoms and score each round's
design against the prior rather than against the realised history, because conditioning on a
history collapses the effective atom count below five by the third round. The obvious
estimator --- the same atoms forming the posterior, choosing the recommendation and grading
it --- is in-sample selection and is badly optimistic: it reports $c = 97\%$ for a design
worth $-64\%$ on held-out atoms. We temper the likelihood, choose the temperature by
leave-one-out over the inference atoms, and grade on a disjoint half.
Appendix~\ref{app:eig} gives the mechanics, the diagnostics, and what each change is worth.
What survives estimation is an ordering, so we report $c$ and $\eta$ as ranks and as means
over dozens of designs, and draw conclusions only from arms compared pairwise at fixed $M$;
Appendix~\ref{app:eig} says what forces that.
